\documentclass[10pt,a4paper]{article}
\usepackage[margin=1.7cm]{geometry}
\usepackage{amsmath,amssymb,amsthm,bm}
\usepackage{booktabs,array,tabularx,ltablex,enumitem,xcolor,hyperref,microtype}
\keepXColumns
\hypersetup{colorlinks,linkcolor=blue!50!black,citecolor=blue!50!black,urlcolor=blue!50!black}
\setlist{nosep,leftmargin=1.2em}
\newcommand{\tr}{\operatorname{tr}}
\newcommand{\Tr}{\operatorname{Tr}}
\newcommand{\ev}[1]{\langle #1\rangle}
\newcommand{\Psib}{\bar\Psi}
\newcommand{\Qb}{\bar Q}
\newcommand{\refnote}[1]{{\small\textcolor{blue!50!black}{[#1]}}}
\title{\vspace{-1.2cm}The quantum-mechanical matrix bootstrap: a cribsheet\\[2pt]
\large with the specialisations needed for purely fermionic (matrix SYK) models}
\author{Matrix SYK bootstrap project}
\date{15 September 2026}
\begin{document}
\maketitle
\vspace{-0.9cm}
\noindent\textit{Scope.} Everything here is drawn from the project library: Lin--Zheng \cite{LZ24,LZ25}, Lin's TASI lectures \cite{TASI}, Laliberte--McPeak \cite{LM}, Klebanov et al.\ \cite{KMPT}, Chen \cite{Chen}, Chang et al.\ \cite{CCSY}. Equation numbers in brackets refer to those papers. Section~\ref{sec:ferm} contains project-specific statements (marked $\star$) that are derived in \texttt{docs/derivations.md} or verified numerically in this repository, not taken from the literature.

\section{Objects}
\begin{description}[style=nextline,font=\bfseries]
\item[State functional.] A linear functional $\phi(O)=\ev{O}$ on operators, meant to be $\langle E|O|E\rangle$ for an energy eigenstate (or $\Tr\rho O$ for a stationary $\rho$, $[\rho,H]=0$) \refnote{TASI (100); LM (1.4)--(1.5)}. Normalisation $\phi(\mathbf 1)=1$ (or $N$). The bootstrap never constructs the state; it constrains $\phi$.
\item[Letters and words.] Letters are the elementary operators ($X_I,P_I,\psi_\alpha$ for BFSS; $\Psi_{ij},\Psib_{ij}$ for us). Words are ordered products. At large $N$ the variables are \emph{single-trace} words $\tr(O_1\cdots O_n)$ with global-symmetry indices contracted into singlets \refnote{LZ24 (8); LZ25 eq.\ after (3)}, normalised $\tr=\frac1N\Tr$, $X\to X/\sqrt N$, so that $\ev{\tr O}=O(1)$ \refnote{LZ24 (7); LZ25 (3)}.
\item[Level (hierarchy).] $\ell(X)=1$, $\ell(P)=2$, $\ell(\psi)=\tfrac32$ (for canonical $\{\psi,\psi\}=\delta$ fermions), additive over letters; chosen so that $\ell([H,O])=\ell(O)+1$ and $\ell(\{Q,O\})=\ell(O)+\tfrac12$ \refnote{LZ24 \S2.5, (25); LZ25 ``Hierarchy''; LM \S3.2.1}. A ``level-$L$ bootstrap'' uses positivity matrices built from operators of level $\le L/2$ (so entries have level $\le L$) and all equality constraints among variables of level $\le L$. ``Free variables'' = variables left after solving the equalities (LZ24 Table~1; LZ25 Table~I).
\end{description}

\section{The constraint catalogue}
\renewcommand{\arraystretch}{1.25}
\noindent\begin{tabularx}{\textwidth}{@{}>{\bfseries}p{2.9cm}Xp{3.6cm}@{}}
\toprule Constraint & Statement & Source \\ \midrule \endfirsthead
\toprule Constraint (cont.) & Statement & Source \\ \midrule \endhead
\bottomrule\endlastfoot
Inner-product positivity &
$M_{ij}=\ev{O_i^\dagger O_j}\succeq 0$ for any finite set $\{O_i\}$: the Gram matrix of $\{O_i|E\rangle\}$. Uses only Hilbert-space positivity, so fermions and sign problems are irrelevant. Adding operators can only tighten bounds (principal minors). &
LM (1.1)--(1.4); LZ25 (5); TASI (95)--(96) \\
Stationarity (EOM) &
$\ev{[H,O]}=0$ for all $O$ (true in any eigenstate or stationary $\rho$). For matrix QM take $O$ single-trace; $[H,\tr O]$ is again single-trace. &
LM (1.5); LZ25 (4); TASI (100) \\
Eigenstate condition &
$\ev{OH}=\ev{HO}=E\ev{O}$: stronger, resolves individual eigenvalues (``archipelago''), but bilinear in $(E,\phi)$ hence non-convex; usable by scanning $E$. Becomes trivial under large-$N$ factorisation. &
LM (2.14) \& fn.\ 2; LZ24 (57)--(58), Fig.\ 4; TASI \S6.2 \\
Reality / time reversal &
$\ev{O_1\cdots O_n}^*=\ev{O_n^\dagger\cdots O_1^\dagger}$; with real wavefunctions $\ev{\tr O_1\cdots O_n}=\pm\ev{\tr O_n\cdots O_1}$ ($-$ for an odd number of $P$'s). Halves the variables. &
LM (3.34); LZ24 \S2.2(2); LZ25 ``Constraints'' \\
Algebra / canonical form &
(Anti)commutation relations reduce every word to a canonical (e.g.\ normal-ordered) form; for matrices, cyclicity of $\tr$ produces \emph{double traces}: $\tr O_1\cdots O_n=\pm\tr O_2\cdots O_nO_1+\text{double trace}$. &
LZ24 (11)--(12), (29); TASI (101) \\
Large-$N$ factorisation &
$\ev{\tr A\,\tr B}\to\ev{\tr A}\ev{\tr B}$ (the \emph{only} place $N=\infty$ enters in the Hamiltonian approach). Makes the constraints quadratic in the variables. Justified by planar dominance; can fail for operators sensitive to flat directions. &
LZ24 (46) \& \S4; TASI \S6.1--6.2 \\
Gauge constraint &
Singlet state: $\ev{\tr(O\,G)}=0$ for all $O$, $G$ the $SU(N)$ generator (e.g.\ $G=i[X,P]-\Psi\Psi^\dagger-\Psi^\dagger\Psi+2N$ for one bosonic + one fermionic matrix). Removes an enormous number of variables; without it the Marinari--Parisi SDP was intractable. &
LZ24 (13); LZ25 after (1); LM (3.12)--(3.13), \S4 \\
Global symmetry (Ward identities, blocking) &
If the state is $G$-invariant, $\ev{O_R}=0$ unless $R$ is the singlet, and $\ev{\bar O_{\bar R\bar r}O_{R r}}=\delta_{\bar RR}\delta_{\bar rr}\,a_{\bar R,R}$. Positivity of the big matrix reduces to $a_{\bar R,R}\succeq0$ block by block; channels are related by crossing kernels ($6j$ symbols). Conserved charges give $\ev{[F,O]}=0$ and block-diagonalise $M$ by charge. &
LZ24 (16)--(24), (32)--(35); LZ25 App.\ F; LM (3.18) \& \S3.2.1 \\
Ground-state positivity &
$N_{ij}=\ev{O_i^\dagger[H,O_j]}\succeq0$: any perturbation $\sum c_iO_i|\Omega\rangle$ has energy $\ge E_0$. Zero-temperature limit of the energy--entropy balance inequality $\ev{O^\dagger O}\log\frac{\ev{O^\dagger O}}{\ev{OO^\dagger}}\le\beta\ev{O^\dagger[H,O]}$. Supplies \emph{upper} bounds and shrinks peninsulas to islands. &
LZ25 (6); TASI (102)--(104); LM (3.20) \\
Supercharge (SUSY vacuum) &
If $Q|\Omega\rangle=0$: $\ev{QO}=\ev{OQ}=0$, at large $N$ $\ev{\{Q_\alpha,O_\alpha\}}=0$ for fermionic $O$. Implies ground-state positivity: $\ev{O^\dagger[H,O]}=\tfrac18\ev{[Q,O]^\dagger[Q,O]}\ge0$. The weaker $\ev{[Q^\dagger Q,O]}=0$ was found redundant with EOM. &
LZ24 (14), App.\ A (47)--(54); TASI (128); LM (3.19) \\
Nonlinear relaxation &
Quadratic terms $x_ix_j$ (from factorisation) $\to$ new variables $\mathcal Q_{ij}$ with $\begin{pmatrix}1&\vec x^{\,T}\\ \vec x&\mathcal Q\end{pmatrix}\succeq0$ and $M\succeq\mathcal Q$ (covariance positivity). Rigorous, convex, weaker than the equality. &
LZ25 App.\ B (B6)--(B7); TASI (98)--(99) \\
\end{tabularx}

\medskip\noindent\textbf{Objective.} Minimise (maximise) $\phi(H)$ for a rigorous lower (upper) bound on $E_0$ (top of spectrum); or fix $E$ and extremise a one-point function to map the allowed region (``peninsula'' if unbounded in some direction, ``island'' if compact). Typical plots: $E$ vs $\ev{\tr X^2}$ (LZ25 Figs.~1--2; LZ24 Figs.~1--2).

\section{Mechanics of a run}
\begin{enumerate}
\item Enumerate variables (single-trace singlet words up to level $L$; or, at finite $N$, all canonical words).
\item Generate and \emph{symbolically solve} the linear equalities (algebra, reality, gauge, Ward, EOM, supercharge) for a set of free variables. This step, not the SDP, dominates cost (LZ24 \S3; LM \S3.2.2: 8\,h of algebra vs 5--10\,min per SDPB point).
\item Build the positivity matrices ($M$, $N$, relaxation blocks) as linear matrix inequalities in the free variables, one block per (charge, irrep).
\item Solve the SDP. Generic solvers (Mathematica, SCS) suffice at low level; the problems lack strict feasibility, so high level needs arbitrary precision: SDPA-GMP (LZ25 App.~A) or SDPB (LM). Bounds should tighten monotonically with level; plateaus followed by jumps happen (LM Fig.~13 and \S3.2.2).
\item Sanity checks: (a) feed the exact state (ED or known solution) into every constraint and verify $M\succeq0$, EOM residual $0$; (b) reproduce analytic low-level bounds, e.g.\ the virial/Polchinski bounds (TASI (114)--(127); LZ25 (7), App.~E).
\end{enumerate}

\section{Worked micro-examples}
\textbf{(a) Fermion number and $E\ge0$ from a $2\times2$ block} (LM (3.26)--(3.28)): with $\{\Psi,\Psi^\dagger\}$ operators,
$\begin{pmatrix}\ev{\tr\Psi^\dagger\Psi}&\ev{\tr\Psi^\dagger\Psi^\dagger}\\ \ev{\tr\Psi\Psi}&\ev{\tr\Psi\Psi^\dagger}\end{pmatrix}\succeq0$ gives $0\le F\le N^2$; combined with the bosonic block it yields $\ev H\ge0$ with $E=0$ only at $F=0$ for the quadratic superpotential.

\smallskip\noindent\textbf{(b) Level-5 bosonic Yang--Mills} (LZ25 App.~E): 19 variables, equalities (E4)--(E32) leave 3 free; irrep blocks (E35)--(E38) give the analytic peninsula/island (E39).

\smallskip\noindent\textbf{(c) SUSY sum of squares} ($\star$, this project): for $H=\{Q,\Qb\}$ include $Q,\Qb$ among the $O_i$. Then $\phi(\Qb Q)$ and $\phi(Q\Qb)$ are diagonal entries of $M\succeq0$, so $\phi(H)\ge0$ is \emph{automatic}; the true content of a lower bound in a lifted sector is whatever forces $\phi(\Qb Q)+\phi(Q\Qb)>0$.

\smallskip\noindent\textbf{(d) Representation-resolved bounds by hand} (KMPT \S3.1--3.2): $C_2\pm\tfrac12A\!\cdot\!A=\tfrac14(A\pm A)^2\ge0$ and Cauchy--Schwarz with a symmetry-averaged density matrix $\rho_R=\frac1{\dim R}\sum_i|e_i\rangle\langle e_i|$ give $E$-vs-Casimir bounds (3.8), (3.21) that are saturated in Casimir-type fermionic matrix models. This is a two-operator bootstrap in each irrep.

\section{Specialisation to purely fermionic matrix models}\label{sec:ferm}
Setting: $p$ complex fermion matrices $\Psi^{(f)}_{ij}$, $\{\Psi_{ij},\Psib_{kl}\}=\delta_{il}\delta_{jk}$, $\Psib_{ij}=\Psi^\dagger_{ji}$, $Q$ cubic in $\Psi$, $H=\{Q,\Qb\}$, R-charge $N_\Psi=\Tr\Psi\Psib$, $[N_\Psi,Q]=3Q$ \refnote{Chen (1.1)--(1.2), (2.4)}. Hilbert space $2^{pN^2}$-dimensional.
\begin{itemize}
\item \textbf{Pauli truncation.} Each mode appears at most once, so the algebra of words closes at finite length; at fixed $N$ the constraint set is finite and the bootstrap can be \emph{exact} at the top level. No flat directions, no continuum of moments.
\item \textbf{Blocks.} $M$ is block-diagonal in $N_\Psi$ (as in LM \S3.2.1: $F=0,\pm1,\pm2$ blocks) and in $SU(N)\times$flavor irrep. Fix a sector with $\phi(N_\Psi)=k$, $\phi(N_\Psi^2)=k^2$: since $\phi((N_\Psi-k)^2)=0$, positivity forces $\phi((N_\Psi-k)O)=0$ for every $O$ in the set. Extremising $\phi(H)$ then bounds the sector ground energy $E_0(k)$ ($\star$).
\item \textbf{Which sector is the ground state.} For $Q=\Tr\Psi^3$, $H=3N(N^2-1)-9\hat C_2$ \refnote{Chen (2.8)}: singlets are the \emph{top} of the spectrum and the BPS states form the maximal-Casimir irrep $r_*$ \refnote{Chen (2.18)--(2.22)}. A gauged (singlet) bootstrap bounds the wrong end. Work ungauged, with irrep-resolved blocks; the singlet-state justification for using open-index operators (LZ25 after (6)) must be replaced by the irrep-averaged functional $\rho_R$ of KMPT (3.15).
\item \textbf{SUSY floor and how to lift it} ($\star$). With $Q,\Qb$ in the basis, $\phi(H)\ge0$ is free. In a lifted sector the SDP must additionally learn that annihilators act first on the sector ground state, e.g.\ $\phi(Q\Qb)=0$ at $N_\Psi=0$. In a basis of elementary letters this follows from $\phi(N_\Psi)=0$, $\phi(\Psi_{ij}\Psib_{ji})\ge0$ and Cauchy--Schwarz. In a basis of \emph{gauge-invariant single traces only} it does not: at $N=3$, levels 3 and 4, all lifted sectors of the single-matrix model return the trivial bound $0$ (verified 2026-09-15). Open-index (adjoint-valued) operators are required in the positivity matrix.
\item \textbf{Ground-state positivity in a sector.} $\ev{O^\dagger[H,O]}\ge0$ is valid for the sector ground state only if $O$ preserves the sector ($[O,N_\Psi]=0$, and likewise for the irrep); a charge-changing $O$ may lower the energy by leaving the sector. For a genuine SUSY vacuum ($E=0$) it is implied by the supercharge constraints (LZ24 App.~A) and adds nothing.
\item \textbf{Concentration as a bootstrap target} \refnote{CCSY Conj.\ 1, \S2; Chen (2.22)}. R-charge concentration = one BPS degree per irreducible complex (labels: $N_\Psi$ mod 3, $SU(N)$ irrep, flavor). A rigorous $E_0(k,\mathbf r)>0$ outside a window of $\le3$ consecutive $k$ proves it; the single-matrix baseline has $E_0=0$ on $N+1$ consecutive sectors.
\item \textbf{Exact single-trace $H$ for $Q=\Tr\Psi^3$} ($\star$, \texttt{docs/derivations.md} D1): $H=\tfrac92\big(\Tr[\Psi\Psi\Psib\Psib]+\Tr[\Psib\Psib\Psi\Psi]\big)-\tfrac32N(N^2-1)$, equivalently $\hat C_2=N N_\Psi-\Tr\Psi\,\Tr\Psib-\Tr[\Psi\Psi\Psib\Psib]$. (An earlier fitted constant $-9(N-1)^2$ is wrong for $N\ge4$.)
\end{itemize}

\section{Numbers to calibrate against}
\begin{itemize}
\item Bosonic YM matrix QM (LZ25 Table~II): $D{=}2$, $M{=}0$: $E\in[0.707832,0.707868]$ at level 14; $D{=}9$: $E\in[6.69946,6.69968]$, $\ev{\tr X_IX_I}\in[2.29195,2.29218]$ at level 11 (Monte Carlo $6.695(5)$, $2.291(1)$). Islands first appear at level 10 for $M=0$.
\item BFSS (LZ24 Table~2): $\ev{\tr X^2}\ge0.355$ at level 9 (Monte Carlo $\approx0.378$); analytic level 6: $\tfrac34(3/50)^{1/3}\approx0.2936$.
\item Marinari--Parisi (LM \S3.2.2): level 8, $44\times44$, 27 free variables; $E_0/N^2\approx0.196\,g^{2/3}$ (WKB $0.242$); no upper bound obtained.
\item Single-matrix fermionic model (Chen; this project): $E\in\{0,18\}$ ($N{=}2$), $\{0,18,45,72\}$ ($N{=}3$); $C_2^{\max}=N(N^2-1)/3$; $\dim r_*=3^{N(N-1)/2}$; BPS window $[N(N-1)/2,N(N+1)/2]$. Three-matrix model at $N=2$: 972 BPS states at $N_\Psi\in\{5,6,7\}$, counts $243{:}486{:}243$, 168 distinct energies (ED, reproduced 2026-09-15).
\end{itemize}

\begin{thebibliography}{9}\small
\bibitem{LZ24} H.~W.~Lin, Z.~Zheng, \emph{Bootstrapping Ground State Correlators in Matrix Theory, Part I}, arXiv:2410.14647.
\bibitem{LZ25} H.~W.~Lin, Z.~Zheng, \emph{High-Precision Bootstrap of Multimatrix Quantum Mechanics}, arXiv:2507.21007.
\bibitem{TASI} H.~W.~Lin, \emph{TASI lectures on Matrix Theory from a modern viewpoint}, arXiv:2508.20970, \S6.
\bibitem{LM} S.~Laliberte, B.~McPeak, \emph{Bootstrapping supersymmetric (matrix) quantum mechanics}, arXiv:2510.01356.
\bibitem{KMPT} I.~R.~Klebanov, A.~Milekhin, F.~Popov, G.~Tarnopolsky, \emph{Spectra of eigenstates in fermionic tensor quantum mechanics}, Phys.\ Rev.\ D 97 (2018) 106023, arXiv:1802.10263.
\bibitem{Chen} Y.~Chen, \emph{Fortuity with a Single Matrix}, arXiv:2511.00790.
\bibitem{CCSY} C.-M.~Chang, Y.~Chen, B.~S.~Sia, Z.~Yang, \emph{Fortuity in SYK models}, JHEP 08 (2025) 003, arXiv:2412.06902.
\end{thebibliography}
\end{document}
